\documentclass[11pt]{amsart}

\usepackage{amsmath,amssymb,amsthm}
\usepackage[margin=1.2in]{geometry}
\usepackage{booktabs}
\usepackage{adjustbox}
\usepackage{xcolor}
\usepackage[colorlinks=true,linkcolor=blue,citecolor=blue,urlcolor=blue]{hyperref}

% Last-resort stretch so a paragraph that cannot break cleanly sets loose
% rather than running into the right margin.
\setlength{\emergencystretch}{2em}

\newtheorem{theorem}{Theorem}[section]
\newtheorem{proposition}[theorem]{Proposition}
\newtheorem{lemma}[theorem]{Lemma}
\newtheorem{corollary}[theorem]{Corollary}
\theoremstyle{definition}
\newtheorem{definition}[theorem]{Definition}
\newtheorem{example}[theorem]{Example}
\newtheorem{question}[theorem]{Question}
\theoremstyle{remark}
\newtheorem{remark}[theorem]{Remark}

\newcommand{\ZZ}{\mathbb{Z}}
\newcommand{\RR}{\mathbb{R}}
\newcommand{\CC}{\mathbb{C}}
\newcommand{\PP}{\mathbb{P}}
\newcommand{\Def}{\operatorname{Def}}
\newcommand{\Sing}{\operatorname{Sing}}
\newcommand{\Spec}{\operatorname{Spec}}
\newcommand{\conv}{\operatorname{conv}}
\newcommand{\FF}{\mathbb{F}}

\title{Non-smoothable Calabi--Yau threefolds \\from reflexive polytopes}
\author{Bernd J. Wuebben}
\subjclass[2020]{14J32 (primary); 14B07, 14M25, 14J17, 52B20 (secondary)}
\keywords{Calabi--Yau threefolds, smoothing of singularities, toric
Gorenstein singularities, versal deformations, reflexive polytopes,
Minkowski decompositions}
\date{July 14, 2026}

\begin{document}

\begin{abstract}
We give explicit examples of compact Calabi--Yau threefolds with isolated
Gorenstein canonical singularities that admit no global smoothing.  The
examples are generic anticanonical hypersurfaces in Gorenstein Fano toric
fourfolds; their singular points are cones over two-dimensional faces of the
corresponding reflexive $4$-polytopes, and smoothability of such cones is
governed by Altmann's description of the versal deformation together with the
smoothing criterion of Corti--Filip--Petracci: the cone over a lattice polygon
with primitive edges is smoothable if and only if the polygon decomposes as a
Minkowski sum of unit segments and standard triangles.  Our first example has
exactly three singular points, each isomorphic to the anticanonical cone over
the Hirzebruch surface $\FF_1$ --- the exceptional local model in Gross's
classification of primitive type~II contractions.  Our second example has
exactly three singular points whose miniversal deformation spaces are
positive-dimensional, so every singular point deforms nontrivially, and yet
the threefold admits no smoothing; to our knowledge this phenomenon has not
previously been exhibited on a compact Calabi--Yau threefold.  A refinement
of the construction, using a third completing vertex to force the relevant
dual edge to have lattice length one, produces a Calabi--Yau threefold whose
singular locus is a \emph{single} point, the anticanonical cone over
$\FF_1$; and a scan of the Kreuzer--Skarke classification produces one whose
singular locus is a single point of the deformable type.  Along the way
we record a complete classification of the isolated Gorenstein toric threefold
singularities whose polygon admits an edge representation with coordinates in
$\{-2,\dots,2\}$ into rigid, deformable-but-non-smoothable, and
smoothable classes, together with the number of smoothing components in each
case.
\end{abstract}

\maketitle
\enlargethispage{5pt}

\section{Introduction}\label{sec:intro}

Let $X$ be a projective threefold with trivial dualizing sheaf,
$h^1(\mathcal{O}_X)=0$, and Gorenstein canonical singularities --- a
\emph{Calabi--Yau threefold} with canonical singularities.  A basic question,
going back to the foundational work of Friedman \cite{Friedman86}, is whether
$X$ can be \emph{smoothed}: does there exist a flat proper family over a
pointed disc with special fibre $X$ and smooth general fibre?

For the mildest singularities the question is completely understood.  If $X$
has ordinary double points, Friedman \cite{Friedman86} showed that a smoothing
exists if and only if the exceptional curves of a small resolution satisfy a
non-trivial ``good'' relation in homology; if $X$ has arbitrary isolated
terminal Gorenstein singularities, Namikawa and Steenbrink
\cite{NamikawaSteenbrink95} proved that $X$ is always smoothable.  Beyond the
terminal case the picture changes qualitatively.  Gross \cite{Gross97}
analyzed the case where $X$ arises from a primitive type~II contraction of a
smooth Calabi--Yau threefold (a divisor $E$ contracted to a point) and proved
that $X$ is then smoothable \emph{unless} $E \cong \PP^2$ or $E \cong \FF_1$
\cite[Theorem~5.8]{Gross97}.  The two exceptions are genuinely
non-smoothable: the germ at the singular point is the anticanonical cone over
$E$, which for $E=\PP^2$ is the rigid quotient singularity
$\frac13(1,1,1)$ \cite{Schlessinger71}, and for $E = \FF_1$ has
one-dimensional tangent space $T^1$ but no genuine deformations at all.
Deformations of Calabi--Yau threefolds with canonical and log canonical
singularities remain an active subject; see Friedman--Laza
\cite{FriedmanLaza25,FriedmanLazaSigma} for recent progress.

This note produces explicit \emph{compact} examples beyond the quotient case,
by a mechanism that is purely combinatorial.  The inputs are:

\begin{enumerate}
\item Altmann's theorem \cite{Altmann97} that for an isolated Gorenstein toric
threefold singularity --- the cone $C(Q)$ over a lattice polygon $Q$ with
primitive edges --- the irreducible components of the reduced miniversal base
correspond to the maximal decompositions of $Q$ into Minkowski sums of lattice
polytopes;
\item the smoothing criterion of Corti--Filip--Petracci
\cite[\S5.1]{CFP}: the \emph{smoothing} components correspond to the Minkowski
decompositions of $Q$ into \emph{unit segments} and \emph{standard triangles};
consequently $C(Q)$ is smoothable if and only if such a decomposition exists;
\item Batyrev's construction \cite{Batyrev94}: for a reflexive polytope
$\Delta \subset \ZZ^4$, the generic anticanonical hypersurface $X$ in the
Gorenstein Fano toric fourfold $\PP_\Delta$ is a Calabi--Yau threefold with
canonical singularities, and (as we spell out in
Proposition~\ref{prop:planting}) the singular points of $X$ are precisely
cones $C(F)$ over the two-dimensional faces $F$ of $\Delta$.
\end{enumerate}

Putting these together: \emph{if a reflexive $4$-polytope $\Delta$ has a
two-dimensional face $F$ whose associated cone $C(F)$ is non-smoothable, then
the generic anticanonical hypersurface $X \subset \PP_\Delta$ is a Calabi--Yau
threefold that admits no smoothing} (Corollary~\ref{cor:global}).  The point
is that non-smoothable faces are easy to recognize (a finite combinatorial
test, Theorem~\ref{thm:trichotomy}) and easy to plant (a small search over
completions of an embedded polygon, \S\ref{sec:examples}).

Our two main examples are the following.  Write $e_1,\dots,e_4$ for the
standard basis of $\ZZ^4$.

\begin{theorem}[= Theorem~\ref{thm:A}]\label{thm:introA}
Let
\[
\Delta_A=\conv\{(0,0,1,0),\,(-2,-1,1,0),\,(-2,-2,1,0),\,(-1,-2,1,0),\,
e_4,\,(1,1,-1,-1)\}\subset\ZZ^4 .
\]
Then $\Delta_A$ is reflexive, and the generic anticanonical hypersurface
$X_A\subset\PP_{\Delta_A}$ is a Calabi--Yau threefold whose singular locus
consists of exactly three points, each analytically isomorphic to the
anticanonical cone over the Hirzebruch surface $\FF_1$.  Every deformation of
$X_A$ over a reduced base preserves these three singular germs; in particular
$X_A$ admits no smoothing.  The maximal projective crepant partial resolution
$\widehat X_A$ is a smooth Calabi--Yau threefold with
$(h^{1,1},h^{2,1})=(5,101)$.
\end{theorem}

The local model here is precisely Gross's exceptional germ: the cone over
$\FF_1$ is \emph{rigid in the reduced sense} --- its miniversal base is a fat
point, with $\dim T^1 = 1$ but no positive-dimensional deformations
\cite[\S5.5]{Gross97}.  Since it is the cone over a non-simplicial polygon, it
is not a quotient singularity, so $X_A$ is not covered by the classical
$\frac13(1,1,1)$ examples.  Compact Calabi--Yau threefolds containing such
points are implicit in the primitive-contraction literature (see
\cite{Gross97,KapustkaKapustka09,Kapustka09}, where the smoothable exceptional
types are systematically smoothed and the $\FF_1$ case is excluded); we are
not aware of an explicit compact example in print, and $\Delta_A$ provides one
lying in the Kreuzer--Skarke classification \cite{KreuzerSkarke00}.

The second example exhibits what we believe is a new phenomenon on compact
Calabi--Yau threefolds.

\begin{theorem}[= Theorem~\ref{thm:B}]\label{thm:introB}
Let
\[
\begin{gathered}
\Delta_B=\conv\{(0,0,1,0),\,(-2,-1,1,0),\,(-3,-2,1,0),\,\\
(-2,-3,1,0),\,(-1,-2,1,0),\,e_4,\,(1,1,-1,-1)\}\subset\ZZ^4 .
\end{gathered}
\]
Then $\Delta_B$ is reflexive, and the generic anticanonical hypersurface
$X_B\subset\PP_{\Delta_B}$ is a Calabi--Yau threefold whose singular locus
consists of exactly three points, each analytically isomorphic to the cone
$C(Q_B)$ over the pentagon $Q_B$ that is the Minkowski sum of the
$\frac13(1,1,1)$-triangle and a unit segment.  Each singular germ has
one-dimensional reduced miniversal base, so each singular point of $X_B$
admits a non-trivial one-parameter local deformation; nevertheless the
miniversal base of each germ has no smoothing component, and $X_B$ admits no
smoothing.  The maximal projective crepant partial resolution
$\widehat X_B$ is a smooth Calabi--Yau threefold with
$(h^{1,1},h^{2,1})=(9,81)$.
\end{theorem}

Thus $X_B$ is a canonical Calabi--Yau threefold all of whose singularities are
honestly deformable --- there is no infinitesimal or local obstruction visible
at first order, and the germs genuinely move --- and yet no deformation ever
smooths it.  The unique local deformation direction pulls the pentagon apart
into its Minkowski summands, trading each pentagon germ for a
$\frac13(1,1,1)$ germ, which is then rigid: the singularity ``cascades'' to
its rigid core and stops.  All previously known non-smoothable canonical
Calabi--Yau examples that we are aware of are locally rigid (quotient points as
in \cite{Gross97}, or the $\FF_1$-cone above).

Both examples have exactly three singular points.  This is not an accident of
the search but a feature of the minimal completion family, and it can be
remedied: a third completing vertex forces the dual edge of the planted face
to have lattice length one, giving a Calabi--Yau threefold whose singular
locus is a \emph{single} non-smoothable point (Theorem~\ref{thm:C}; the same
device produces a Calabi--Yau threefold with a single $\frac13(1,1,1)$ point,
where the classical $X_9$ has three).  The deformable analogue exists as
well: a systematic scan of the Kreuzer--Skarke classification
\cite{KreuzerSkarke00} locates a $10$-vertex reflexive polytope whose
hypersurface has a single singular point of the pentagon type $C(Q_B)$
(Theorem~\ref{thm:D}) --- a compact Calabi--Yau threefold with one singular
point which deforms non-trivially and yet is never absorbed by a smoothing.

The two theorems are instances of a general planting mechanism which we state
as Proposition~\ref{prop:planting} and Corollary~\ref{cor:global}, and which
turns the smoothing question for a large class of compact Calabi--Yau
threefolds into polygon combinatorics.  As bookkeeping for that combinatorics
we record in \S\ref{sec:local} a complete census of the isolated Gorenstein
toric threefold singularities whose polygon has an edge representation with
coordinates in $[-2,2]$: there are $217$ classes up to equivalence --- the
ordinary double point, $8$ rigid classes, $79$ deformable-but-non-smoothable
classes, and $129$ smoothable classes.  The five classes whose polygon has one
interior lattice point are exactly the anticanonical cones over the five
smooth toric del Pezzo surfaces, and the census reproduces the classical
deformation-theoretic data of these cones \cite[\S5.5]{Gross97} exactly
(Remark~\ref{rem:delpezzo}); this both validates the computation and places
the new examples in context: non-smoothability is common among Gorenstein
toric threefold singularities (87 of the 216 canonical classes in the census),
while the del Pezzo cone catalogue sees only its smallest stratum.

Our approach is the Calabi--Yau analogue, one dimension up, of the program of
Petracci and Corti--Hacking--Petracci on (non-)smoothability of Gorenstein
toric \emph{Fano} threefolds \cite{Petracci20,CHP24}, where the local models
are the same cones $C(F)$, attached to the facets of a reflexive
$3$-polytope.  In our setting the local-to-global step is simpler than in
loc.~cit.: because the relevant local models are non-smoothable outright, no
almost-flatness or bundle-cohomology obstruction is needed --- the versal base
of the germ already has no smoothing component, and this kills all global
smoothings (Lemma~\ref{lem:localglobal}).

All computations in this paper (the census, the face inventories of the
example polytopes, the dual-edge point counts, the planting statistics and
the Kreuzer--Skarke scan of \S\ref{sec:questions}, and the Hodge numbers) are
implemented in six short exact-arithmetic computer programs, validated
against the classical anchors recalled in the text (the quintic; $X_9 \subset
\PP(1,1,1,3,3)$; the del Pezzo cone catalogue; Altmann's hexagon; and, for
the scan, the Hodge numbers recorded in the Kreuzer--Skarke data itself).
They are included with the arXiv submission as ancillary files.

\subsection*{Acknowledgements}
I am grateful to Mark Gross for first drawing my attention, many years ago, to the problem of smoothing canonical singularities on Calabi--Yau threefolds.

\section{Local models: cones over polygons}\label{sec:local}

\subsection{The dictionary}
Let $N \cong \ZZ^3$ and let $Q \subset N_\RR$ be a lattice polygon placed at
height one, i.e.\ $Q \subset \{ \langle u^*, \cdot\rangle = 1\}$ for a
primitive $u^* \in M = N^\vee$.  Write $\sigma_Q$ for the cone over $Q$ and
\[
C(Q) \;=\; \Spec \CC[\sigma_Q^\vee \cap M].
\]
Every Gorenstein toric affine threefold without torus factors arises this way,
and $C(Q)$ depends only on $Q$ up to affine unimodular equivalence of the
plane $\{ \langle u^*,\cdot\rangle = 1\}$, i.e.\ up to $\ZZ^2 \rtimes
GL_2(\ZZ)$ acting on $Q$.  We freely identify $Q$ with a polygon in $\ZZ^2$.
The basic dictionary (see e.g.\ \cite{CFP}) is:
\begin{itemize}
\item $C(Q)$ has an \emph{isolated} singularity if and only if every edge of
$Q$ has lattice length one (\emph{unit edges}); an edge of lattice length
$m\ge 2$ produces a curve of transverse $A_{m-1}$-singularities.
\item $C(Q)$ is canonical and Gorenstein; it is \emph{terminal} if and only if
$Q$ has no lattice points other than its vertices, and it is smooth if and
only if $Q$ is a standard triangle (unimodular image of
$\conv\{(0,0),(1,0),(0,1)\}$).
\item For a unit-edge $k$-gon, Pick's theorem gives $i = A - k/2 + 1$ interior
lattice points, where $A$ is the euclidean area.
\end{itemize}
Since translations do not change the edges, a unit-edge polygon is equivalent
data to its multiset $E(Q)=\{e_1,\dots,e_k\}$ of primitive edge vectors,
listed counterclockwise, with $\sum e_i = 0$; and $GL_2(\ZZ)$-equivalence of
polygons is $GL_2(\ZZ)$-equivalence of edge multisets.

\subsection{Deformations: the trichotomy}
Minkowski decompositions of unit-edge polygons are transparent at the level of
edge multisets.

\begin{lemma}\label{lem:partition}
Let $Q$ be a unit-edge polygon with edge multiset $E(Q)$.  Minkowski
decompositions $Q = R_1 + \dots + R_r$ into lattice polytopes (up to
translation of the summands) correspond bijectively to partitions
$E(Q)=E_1\sqcup\dots\sqcup E_r$ into non-empty zero-sum sub-multisets.  Under
this correspondence a summand is a unit segment if and only if its part is a
pair $\{v,-v\}$, and a standard triangle if and only if its part is a triple
$\{a,b,c\}$ with $a+b+c=0$ and $|\det(a,b)|=1$.
\end{lemma}

\begin{proof}
If $Q = R + S$, then each edge of $Q$ decomposes as the sum of the faces of
$R$ and $S$ with the same outer normal, and lattice lengths add; since the
edges of $Q$ are primitive, each edge comes entirely from $R$ or from $S$.
Hence $E(Q) = E(R) \sqcup E(S)$, and inductively any decomposition partitions
$E(Q)$ into zero-sum parts.  Conversely, a zero-sum sub-multiset of primitive
vectors, sorted counterclockwise, closes up to a convex lattice polygon
(a segment when the part is $\{v,-v\}$), and the Minkowski sum of the polygons
associated to the parts of a partition has edge multiset $E(Q)$; since a
convex polygon is determined by its edge multiset up to translation, that sum
is $Q$.  The last statement is immediate: a triangle with primitive edges
$a,b,c=-a-b$ is unimodular if and only if $|\det(a,b)|=1$.
\end{proof}

The deformation theory of $C(Q)$ is completely described by two results.

\begin{theorem}[Altmann \cite{Altmann97}; Corti--Filip--Petracci
{\cite[\S5.1]{CFP}}]\label{thm:trichotomy}
Let $Q$ be a unit-edge lattice polygon, not a standard triangle, and let
$V = C(Q)$.
\begin{enumerate}
\item The irreducible components of the reduced miniversal base space of $V$
are in bijection with the maximal Minkowski decompositions of $Q$ into lattice
summands; each component is smooth, and the component associated to a
decomposition with $r$ summands has dimension $r-1$.
\item The smoothing components are exactly those associated to decompositions
all of whose summands are unit segments and standard triangles.
\end{enumerate}
Consequently, with the notation of Lemma~\textup{\ref{lem:partition}}, exactly
one of the following holds:
\begin{itemize}
\item[\textup{(R)}] \textup{(rigid)} $E(Q)$ has no proper non-empty zero-sum
sub-multiset.  Then the reduced miniversal base is a point: $V$ has no
positive-dimensional deformations, and is not smoothable.
\item[\textup{(D)}] \textup{(deformable, non-smoothable)} $E(Q)$ has a proper
zero-sum sub-multiset, but admits no partition into pairs $\{v,-v\}$ and
unimodular triples.  Then $V$ deforms non-trivially but no deformation has
smooth general fibre.
\item[\textup{(S)}] \textup{(smoothable)} $E(Q)$ admits such a partition; the
number of such partitions equals the number of smoothing components of the
miniversal base.
\end{itemize}
\end{theorem}

For the dimensions in (1) and the concentration of $T^1$ in the Gorenstein
degree see also \cite{Altmann95}; for unit-edge $k$-gons one has
$\dim T^1 = k-3$.  Filip \cite{Filip25} re-derives and extends part (2) to
non-isolated Gorenstein toric singularities via $0$-mutable Laurent
polynomials.

\begin{example}\label{ex:pentagon}
The two smallest illustrations of the three classes:
\begin{itemize}
\item The $\frac13(1,1,1)$ singularity $=\CC^3/\mu_3$ is the cone over the
triangle $T$ with $E(T)=\{(-2,-1),(1,-1),(1,2)\}$; no proper subset sums to
zero, so it is rigid --- recovering Schlessinger \cite{Schlessinger71}.
\item The ordinary double point is the cone over the unit square,
$E=\{\pm(1,0),\pm(0,1)\}$: one partition into two segment pairs, hence
smoothable with a single smoothing component ($xy-zw=t$).
\item The pentagon $Q_B = T + s$, the Minkowski sum of $T$ with a unit segment
$s$, has $E(Q_B)=\{(-2,-1),(-1,-1),(1,-1),(1,1),(1,2)\}$.  Its only proper
zero-sum sub-multisets are the pair $\{(-1,-1),(1,1)\}$ and the complementary
triple $E(T)$, which is not unimodular ($\det = 3$).  So $Q_B$ has exactly one
non-trivial maximal decomposition, $T+s$: the reduced miniversal base of
$C(Q_B)$ is a smooth curve, and there is no smoothing component --- $C(Q_B)$
is deformable but not smoothable.  (Here $\dim T^1 = 2$, so the miniversal
base is in addition non-reduced.)
\end{itemize}
\end{example}

\begin{lemma}[The cascade]\label{lem:cascade}
Let $Q = R_0 + \dots + R_r$ be a Minkowski decomposition of a unit-edge
polygon into lattice summands, let $f_0,\dots,f_r$ be the standard
basis of $\ZZ^{r+1}$, and let $\widetilde V$ be the Gorenstein toric
$(3+r)$-fold associated to the \emph{Cayley cone}
$\widetilde\sigma=\operatorname{Cone}\bigl(\bigcup_{i} R_i\times\{f_i\}\bigr)
\subset \RR^2\oplus\RR^{r+1}$; write $u_0,\dots,u_r$ for the dual basis of
$f_0,\dots,f_r$, viewed as characters of $\widetilde V$, and let
\[
\pi=\bigl(\chi^{u_0}-\chi^{u_1},\ \dots,\ \chi^{u_0}-\chi^{u_r}\bigr)\colon
\widetilde V \longrightarrow \CC^r,
\qquad \pi^{-1}(0)\cong C(Q),
\]
be Altmann's homogeneous deformation attached to the decomposition
\cite{Altmann95} \textup{(}cf.\ \cite[Remark~6.2]{CFP}\textup{)}.  Then:
\begin{enumerate}
\item the general fibre of $\pi$ has exactly one singular point of analytic
type $C(R_i)$ for each two-dimensional non-smooth summand $R_i$, and no
other singular points;
\item for the pentagon $Q_B = T + s$ of
Example~\textup{\ref{ex:pentagon}}, the family $\pi$ dominates the
positive-dimensional component of $\Def C(Q_B)$, whose general fibre
therefore has a single singular point, of type $\frac13(1,1,1)$.
\end{enumerate}
\end{lemma}

\begin{proof}
(1)  For two-dimensional $R_i$ the
orbit closure of the face $\operatorname{Cone}(R_i\times\{f_i\})$ is an
$r$-dimensional affine toric variety whose character lattice is freely
generated by the restrictions of the $u_j$, $j \neq i$, while
$\chi^{u_i}$ vanishes on it; so for generic $t\in\CC^r$ the $r$ equations of
the fibre $\pi^{-1}(t)$ determine on it the single reduced point
$\chi^{u_0}=t_i$, $\chi^{u_j}=t_i-t_j$, lying in the open orbit, and the
implicit-function argument in the proof of
Proposition~\ref{prop:planting}(2) below identifies the germ of
$\pi^{-1}(t)$ there with $(C(R_i),0)$.  On the orbit closure of a face of
$\widetilde\sigma$ meeting two distinct summand heights, some equation of
$\pi^{-1}(t)$ restricts to $0=t_i$ or $t_i=t_j$, so the general fibre avoids
these strata.  The remaining strata are the open orbit and the orbits of
the proper faces $\tau$ of the cones
$\operatorname{Cone}(R_i\times\{f_i\})$, and there the fibre is smooth.
Over the open orbit this holds because $u_0,\dots,u_r$ extend to a basis of
the character lattice.  Along the orbit of a proper face $\tau$ the ambient
$\widetilde V$ is itself smooth: $\tau$ is spanned either by a single ray
$(v,f_i)$ with $v$ a vertex of $R_i$, which is primitive, or by the rays
over an edge $[v,v']$ of $R_i$, and since the edge is unit the pair
$(v,f_i),\ (v'-v,0)$ extends to a lattice basis.  The functionals $u_j$,
$j \neq i$ (including $u_0$), vanish on $\tau$ and form part of a basis of
the saturated lattice $\tau^\perp \cap M$, so on the smooth chart
$U_\tau \cong \CC^{\dim\tau}\times(\CC^\ast)^{3+r-\dim\tau}$ the characters
$\chi^{u_0}$ and $\chi^{u_j}$, $j\neq i$, may be taken as coordinates of
the torus factor, while $\chi^{u_i}$ is divisible by a nonconstant monomial
in the $\CC^{\dim\tau}$-coordinates and hence vanishes along the orbit
together with its derivatives in the torus coordinates.  The Jacobian of
the $r$ equations of $\pi^{-1}(t)$ with respect to
$\chi^{u_0},\chi^{u_j}$ is therefore unitriangular, and the fibre is a
smooth complete intersection near the orbit.

(2)  By Example~\ref{ex:pentagon} the reduced miniversal base of $C(Q_B)$
is a smooth curve and $T+s$ is the unique non-trivial maximal
decomposition, with $r=1$.  The family $\pi$ is non-trivial --- by (1) its
general fibre has a single singular point of type
$C(T) = \frac13(1,1,1)$, not $C(Q_B)$ --- so its classifying map to the
miniversal base is non-constant and dominates the curve.  Hence the
general fibre over the positive-dimensional component of $\Def C(Q_B)$ is
the general fibre of $\pi$: a single $\frac13(1,1,1)$ point.  The only
available deformation of the germ thus trades the pentagon cone for the
rigid quotient cone, after which no further deformation is possible.
\end{proof}

\subsection{A census, and the del Pezzo catalogue}\label{subsec:census}
Both tests in Theorem~\ref{thm:trichotomy} are finite, so the trichotomy can
be tabulated.  We enumerated all unit-edge polygons whose (counterclockwise)
edge vectors have coordinates in $[-2,2]$ --- all $k$-gons with
$3 \le k \le 16$ occur --- and reduced modulo $GL_2(\ZZ)$.

\begin{proposition}\label{prop:census}
Up to equivalence there are exactly $217$ unit-edge polygons with an edge
representation in the box above, excluding the standard triangle.  They
comprise: the unit square (the ordinary double point, the unique terminal
class, smoothable); $8$ rigid classes; $79$ deformable-but-non-smoothable
classes; and $129$ smoothable classes.  The $8$ rigid classes are, listed as
$[k,i;\,E(Q)]$:
\[
\begin{array}{l}
[3,1;\,(-2,-1),(1,-1),(1,2)]\\[2pt]
[4,1;\,(-2,-1),(0,-1),(1,0),(1,2)]\\[2pt]
[4,2;\,(-2,-1),(2,-1),(1,2),(-1,0)]\\[2pt]
[4,2;\,(-2,-1),(2,-1),(1,1),(-1,1)]\\[2pt]
[5,2;\,(-2,-1),(2,-1),(1,1),(0,1),(-1,0)]\\[2pt]
[5,3;\,(-2,-1),(-1,-1),(2,-1),(1,2),(0,1)]\\[2pt]
[5,4;\,(-2,-1),(0,-1),(2,-1),(1,2),(-1,1)]\\[2pt]
[6,5;\,(-2,-1),(-1,-1),(2,-1),(1,0),(1,2),(-1,1)]
\end{array}
\]
\end{proposition}

The count is computer-verified, and it is worth stating exactly what is
certified.  The per-class verdicts and smoothing-component counts are
$GL_2(\ZZ)$-invariants of the edge multiset, hence correct independently of
any deduplication; they are anchored against the classical cases (the
conifold, the rigid $\frac13(1,1,1)$ cone, Altmann's $\mathrm{dP}_6$ hexagon
with its two smoothing components, and the del~Pezzo catalogue of
Remark~\ref{rem:delpezzo}).  The class count $217$ depends on the
deduplication, which is a union--find with transitive closure over the
$GL_2(\ZZ)$-elements with entries in $[-3,3]$; the count is unchanged when
the entry bound is raised to $[-4,4]$ (both runs are performed by the
ancillary code on every execution), and this stability is the certificate we
offer --- we have not certified the count against a trusted normal form.  Note that rigidity does not persist only
for few edges: for every $k$ there are rigid unit-edge $k$-gons, e.g.\ with
$E(Q) = \{(1,a_1),\dots,(1,a_{k-1}),\, (-(k-1),-\textstyle\sum a_j)\}$ for
distinct $a_j$ with $\gcd(k-1,\sum a_j)=1$: any zero-sum subset must contain
the last vector and then all the others.

\begin{remark}[del Pezzo cross-check]\label{rem:delpezzo}
The census contains exactly five classes with $i=1$ interior point.  If $Q$ is
unit-edge with $i=1$, then $Q$ is reflexive (after centering) and one checks
that $\sigma_Q^\vee$ is the cone at height one over the polar dual $Q^\circ$;
hence
\[
C(Q) \;\cong\; \Spec \textstyle\bigoplus_{m\ge 0} H^0\!\left(S,\,
\omega_S^{-m}\right),
\]
the anticanonical cone over the toric surface $S$ with moment polygon
$Q^\circ$.  For the five census classes the surfaces $S$ are smooth: they are
precisely the five smooth toric del Pezzo surfaces, and the census verdicts
reproduce the classical deformation theory of their cones
\cite[\S5.5]{Gross97} (which rests on \cite{Altmann97}):
\begin{center}
\adjustbox{max width=\textwidth}{%
\begin{tabular}{lcccc}
\toprule
$S$ & $\deg S$ & $[k,i]$ & $\dim T^1 = k-3$ & verdict \\
\midrule
$\PP^2$ & $9$ & $[3,1]$ & $0$ & rigid \\
$\FF_1$ & $8$ & $[4,1]$ & $1$ & rigid (fat point) \\
$\PP^1{\times}\PP^1$ & $8$ & $[4,1]$ & $1$ & smoothable, $1$ component \\
$\mathrm{dP}_7$ & $7$ & $[5,1]$ & $2$ & smoothable, $1$ component \\
$\mathrm{dP}_6$ & $6$ & $[6,1]$ & $3$ & smoothable, $2$ components \\
\bottomrule
\end{tabular}}
\end{center}
In particular the second rigid class of Proposition~\ref{prop:census},
$Q_A$ with $E(Q_A)=\{(-2,-1),\allowbreak(0,-1),\allowbreak(1,0),\allowbreak(1,2)\}$, satisfies
$Q_A^\circ = $ the moment polygon of $(\FF_1, -K)$, so
\emph{$C(Q_A)$ is the anticanonical cone over $\FF_1$} --- the exceptional
germ of \cite[Theorem~5.8]{Gross97}.  The deformable-but-non-smoothable
classes all have $i \ge 2$; they are cones over polarized toric surfaces that
are not anticanonically polarized del Pezzo surfaces, and so lie outside the
catalogue above.
\end{remark}

\section{Two-faces of reflexive polytopes and Batyrev hypersurfaces}
\label{sec:planting}

Let $\widetilde N \cong \ZZ^4$ and let $\Delta \subset \widetilde N_\RR$ be a
reflexive polytope: a lattice polytope with $0$ in its interior whose polar
dual $\Delta^\circ=\{u : \langle u, v\rangle \ge -1 \ \forall v \in \Delta\}$
is again a lattice polytope.  Let $\PP_\Delta$ be the toric fourfold of the
face fan of $\Delta$; it is Gorenstein Fano, $-K_{\PP_\Delta}$ is ample and
globally generated, and its space of sections has the monomial basis
$\{\chi^m : m \in \Delta^\circ \cap \widetilde M\}$ \cite{Batyrev94,CLS}.

For a face $G \preceq \Delta$ write $\sigma_G$ for the cone over $G$,
$V(\sigma_G) \subset \PP_\Delta$ for the corresponding orbit closure, and
$G^\circ \preceq \Delta^\circ$ for the dual face.  If $F$ is a
two-dimensional face, then $\sigma_F$ is a three-dimensional cone,
$V(\sigma_F)$ is a toric curve $\cong \PP^1$, and $F^\circ$ is an edge of
$\Delta^\circ$; we write $\ell(F^\circ)$ for its lattice length.

\emph{Normalization convention.}  When listing example data we describe each
facet of $\Delta$ by the primitive integral functional $u$ normalized so that
the facet lies on $\{\langle u,\cdot\rangle = 1\}$; the corresponding vertex
of $\Delta^\circ$ is then $-u$, and the dual face $G^\circ$ of a face $G$ is
spanned by the \emph{negatives} of the facet functionals through $G$.  Since
lattice length is invariant under $u \mapsto -u$, we compute $\ell(F^\circ)$
directly on the listed functionals.  (This matches the normalization used by
the ancillary programs.)  The
supporting functionals of the two facets of $\Delta$ containing $F$ restrict
to the sublattice $N_F = \widetilde N \cap \RR\sigma_F$ as an integral height
function equal to $-1$ on $F$; thus $(\sigma_F, N_F)$ is the Gorenstein cone
over the lattice polygon $F$, and we write $C(F)$ for the associated toric
threefold as in \S\ref{sec:local}, computed in the lattice induced on the
affine span of $F$.

\begin{proposition}\label{prop:planting}
Let $\Delta \subset \ZZ^4$ be reflexive, all of whose edges have lattice
length one, and let $X \in |{-K_{\PP_\Delta}}|$ be generic.  Then:
\begin{enumerate}
\item $X$ is a Calabi--Yau threefold with Gorenstein canonical singularities.
\item For each two-dimensional face $F \preceq \Delta$ with $C(F)$ singular,
$X$ meets the curve $V(\sigma_F)$ in exactly $\ell(F^\circ)$ points, all in
the open orbit, and at each such point $p$ the germ $(X,p)$ is analytically
isomorphic to $(C(F), 0)$.
\item $X$ is smooth away from these points.
\end{enumerate}
\end{proposition}

\begin{proof}
(1) is Batyrev's theorem \cite[\S4]{Batyrev94}.

(2) Fix such an $F$.  The restriction map on anticanonical sections sends
$\chi^m$ to a non-zero section on $V(\sigma_F)$ exactly when
$m \in F^\circ \cap \widetilde M$: indeed $\chi^m$ vanishes identically on
$V(\sigma_F)$ unless $\langle m, \cdot \rangle$ attains its minimum $-1$ on
$F$.  The restricted line bundle is the toric line bundle on
$V(\sigma_F)\cong\PP^1$ with one-dimensional moment polytope $F^\circ$, of
degree $\ell(F^\circ)$ \cite[\S6.3]{CLS}, and the lattice points of $F^\circ$
restrict to the full monomial basis of $H^0(\PP^1,
\mathcal{O}(\ell(F^\circ)))$.  Hence for generic $X$ the intersection
$X \cap V(\sigma_F)$ consists of $\ell(F^\circ)$ distinct reduced points
avoiding the two torus-fixed points.

For the local structure, split $\widetilde N \cong N_F \oplus \ZZ$; then the
open toric chart of $\sigma_F$ is $U_{\sigma_F} \cong C(F) \times \CC^\ast$,
with $V(\sigma_F)\cap U_{\sigma_F} = \{0\} \times \CC^\ast$
\cite[Prop.~3.3.9]{CLS}.  Trivialize $-K$ on this chart and embed
$C(F) \subset \CC^r$ by a generating set of $\sigma_F^\vee \cap M_F$; the
defining equation of $X$ extends to a function $\widetilde f(z, t)$ on
$\CC^r\times\CC^\ast$.  At a point $p=(0,t_0)$ of the transverse intersection
above, $\widetilde f(0, \cdot)$ has a simple zero at $t_0$, so
$\partial\widetilde f/\partial t (p) \neq 0$, and by the implicit function
theorem $\{\widetilde f = 0\}$ is, near $p$, the graph of an analytic function
$t = \varphi(z)$.  Intersecting with $C(F)\times\CC^\ast$ and projecting along
the graph yields an analytic isomorphism $(X, p) \cong (C(F), 0)$.

(3) $\Sing \PP_\Delta$ is the union of the orbit closures of the non-smooth
cones of the face fan.  Cones over vertices and the full-dimensional cones'
smooth loci are irrelevant for a generic hypersurface: since
$-K_{\PP_\Delta}$ is globally generated, Bertini gives that $X$ is smooth away
from $\Sing\PP_\Delta$, and $X$ avoids the finitely many torus-fixed points
(the section $\chi^{u_\Gamma}$, for $u_\Gamma$ the vertex of $\Delta^\circ$
dual to a facet $\Gamma$, is non-vanishing at the fixed point of
$\sigma_\Gamma$).  For a two-dimensional cone $\sigma_e$ over an edge $e$ of
$\Delta$ with vertices $v, w$: since $\ell(e)=1$ and $v$ lies at height $1$
with respect to the (restriction of the) supporting functional of any facet
containing $e$, the pair $(v, w-v)$ is a basis of $\widetilde N \cap
\RR\sigma_e$, so $\sigma_e$ is smooth.  Hence, under the all-unit-edges
hypothesis, $\Sing\PP_\Delta$ meets $X$ only in the curves $V(\sigma_F)$ of
part (2), and there only at the listed points.
\end{proof}

\begin{remark}
If $\Delta$ has an edge $e$ of lattice length $m \ge 2$, the same analysis
shows that $X$ acquires a curve of transverse $A_{m-1}$-singularities along
$V(\sigma_e)$; we exclude this only to keep the singular locus zero-dimensional.
\end{remark}

\section{Local-to-global: no smoothing}\label{sec:nosmoothing}

\begin{lemma}\label{lem:localglobal}
Let $X$ be a reduced projective variety with an isolated singular point $p$
such that no irreducible component of the reduced miniversal base space of the
germ $(X,p)$ has smooth general fibre.  Then for every flat proper family
$\mathcal{X} \to (T,0)$ over a reduced irreducible pointed base with
$\mathcal{X}_0 \cong X$, every fibre $\mathcal{X}_t$ for $t$ near $0$ is
singular.  In particular $X$ admits no smoothing.

If moreover the reduced miniversal base of $(X,p)$ is a point, then the
induced family of germs at $p$ is trivial: every $\mathcal{X}_t$ has a
singular point with germ isomorphic to $(X,p)$.
\end{lemma}

\begin{proof}
By Grauert \cite{Grauert72}, the isolated singularity $(X,p)$ has a miniversal
deformation $\mathcal{S} \to (S,0)$, which we represent on a suitable Milnor
representative; versality holds for analytic families over reduced base germs
after shrinking \cite[II.1]{GLS}.  The relative singular locus
$\Sigma \subset \mathcal{S}$ is closed analytic, and (after shrinking) the
projection $\Sigma \to S$ is finite, so its image is a closed analytic subset
$D \subseteq S$ --- the locus of singular fibres.  A component of $S$ is a
smoothing component precisely when it is not contained in $D$; by hypothesis
there are none, so $D = S$ and \emph{every} fibre of the miniversal family is
singular.

Now let $\mathcal{X} \to (T,0)$ be as in the statement and restrict it to a
Milnor representative around $p$; after shrinking $T$ this restriction is
induced from $\mathcal{S} \to S$ by a base-change map $h : (T,0) \to (S,0)$.
Its fibres are fibres of the miniversal family, hence singular; so
$\mathcal{X}_t$ has a singular point near $p$ for every $t$ near $0$.

For the final statement: if $S_{\mathrm{red}}$ is a point, then $h$ factors
through $S_{\mathrm{red}}$ because $T$ is reduced, i.e.\ $h$ is constant, and
the induced family of germs is the pullback of the central fibre.
\end{proof}

\begin{remark}
Lemma~\ref{lem:localglobal} is deliberately the crudest possible
local-to-global statement: it requires no Calabi--Yau hypothesis and no
obstruction theory, because the local model already forbids smooth nearby
fibres.  For the much finer relationship between local and global
deformations of a Calabi--Yau threefold with isolated canonical
singularities --- in particular when local smoothings do exist and the
question is whether they are realized by global deformations --- see Gross
\cite[\S\S1--2]{Gross97} and Namikawa's stratified local moduli
\cite{Namikawa02}.
\end{remark}

\begin{corollary}\label{cor:global}
Let $\Delta$ be a reflexive $4$-polytope with unit edges possessing a
two-dimensional face $F$ such that the polygon $F$ is of type \textup{(R)} or
\textup{(D)} in Theorem~\textup{\ref{thm:trichotomy}}.  Then the generic
anticanonical hypersurface $X \subset \PP_\Delta$ is a Calabi--Yau threefold
with canonical Gorenstein singularities which admits no smoothing.
\end{corollary}

\begin{proof}
By Proposition~\ref{prop:planting}, $X$ has $\ell(F^\circ) \ge 1$ points with
germ $C(F)$; by Theorem~\ref{thm:trichotomy} the miniversal base of this germ
has no smoothing component; by Lemma~\ref{lem:localglobal} no deformation of
$X$ over a reduced irreducible base (in particular no smoothing over a disc)
has smooth general fibre.
\end{proof}

\section{The examples}\label{sec:examples}

Both examples plant a non-smoothable polygon $Q$ from \S\ref{sec:local} as a
two-dimensional face of a reflexive $4$-polytope of the simple shape
$\conv\left(Q \times \{(1,0)\} \cup \{e_4,\ w\}\right)$, found by a small
computer search over the completing vertex $w$.  (Placing $Q$ in the plane
$x_3=1,\,x_4=0$ embeds it in a saturated affine sublattice, so the face
polygon carries the intrinsic lattice structure of $Q$; translations of $Q$
within that plane and the normalization of the first completing vertex to
$e_4$ are unimodular coordinate changes.)

\subsection{A Calabi--Yau threefold with three $\FF_1$-cone points}

\begin{theorem}\label{thm:A}
Let $Q_A$ be the rigid quadrilateral with
$E(Q_A)=\{(-2,-1),\allowbreak(0,-1),\allowbreak(1,0),\allowbreak(1,2)\}$, embedded with vertices
$(0,0),(-2,-1),(-2,-2),(-1,-2)$, and let
\[
\Delta_A=\conv\bigl(Q_A\times\{(1,0)\}\ \cup\ \{e_4,\ (1,1,-1,-1)\bigr\}
\subset \ZZ^4 .
\]
Then:
\begin{enumerate}
\item $\Delta_A$ is reflexive with $6$ vertices and $6$ facets, and all $13$
edges of $\Delta_A$ have lattice length one.  Its $13$ two-dimensional faces
consist of the face $F_A = Q_A\times\{(1,0)\}$ and $12$ standard triangles.
\item The generic anticanonical hypersurface $X_A \subset \PP_{\Delta_A}$ is a
Calabi--Yau threefold whose singular locus consists of exactly
$\ell(F_A^\circ) = 3$ points, at each of which the germ of $X_A$ is the
anticanonical cone over $\FF_1$.
\item Every deformation of $X_A$ over a reduced irreducible base preserves the
three singular germs up to analytic isomorphism.  In particular $X_A$ admits
no smoothing.
\item Batyrev's formulas give $(h^{1,1},h^{2,1}) = (5,101)$, hence Euler
number $-192$, for the maximal projective crepant partial resolution
$\widehat X_A$.
\end{enumerate}
\end{theorem}

\begin{proof}
(1) is a finite verification.  The six facet inequalities
$\langle u, \cdot\rangle \le 1$ of $\Delta_A$ have primitive normals
\[
(0,0,1,1),\ (0,0,1,-2),\ (-3,6,1,1),\ (6,-3,1,1),\ (-3,0,-5,1),\
(0,-3,-5,1),
\]
each at level exactly $1$, which is reflexivity; the face $F_A$ is the
intersection of $\Delta_A$ with the two supporting hyperplanes of the first
two normals, and the classification of the remaining twelve two-dimensional
faces and of all edge lengths is a direct computation (independently verified
in exact arithmetic).  The dual edge $F_A^\circ$ is spanned by the negatives
of the first two normals (the convention of \S\ref{sec:planting}), of lattice
length $\gcd\bigl((0,0,1,1)-(0,0,1,-2)\bigr)=3$.

(2) By Proposition~\ref{prop:planting}, the singular points of $X_A$ are: for
$F_A$, three points with germ $C(Q_A)$; for the twelve standard-triangle
faces, none, as their cones are smooth.  By Remark~\ref{rem:delpezzo},
$C(Q_A)$ is the anticanonical cone over $\FF_1$.

(3) $Q_A$ is rigid: $E(Q_A)$ has no proper zero-sum sub-multiset (a direct
check of the $14$ proper non-empty subsets).  By
Theorem~\ref{thm:trichotomy}(R) the reduced miniversal base of each germ is a
point, and Lemma~\ref{lem:localglobal} applies, including its final statement.

(4) With $\ell(\Delta_A)=8$ and $\ell(\Delta_A^\circ)=130$ lattice points and
the face data above, Batyrev's formulas \cite{Batyrev94} for the Hodge numbers
of $\widehat X_A$ give $h^{1,1}=5$ and $h^{2,1}=101$.  (The same
implementation returns $(1,101)$ for the quintic simplex and $(4,112)$ for
$X_9\subset\PP(1,1,1,3,3)$.)
\end{proof}

\begin{remark}
The three germs of $X_A$ have $\dim T^1 = 1$ but fat-point miniversal bases:
$X_A$ has first-order deformations at each singular point, all obstructed.
This is exactly the local model excluded in Gross's smoothability theorem for
primitive type~II contractions \cite[Theorem~5.8, \S5.5]{Gross97}, here
realized three times on an explicit compact Calabi--Yau threefold in the
Kreuzer--Skarke list.
\end{remark}

\subsection{A deformable but non-smoothable Calabi--Yau threefold}

\begin{theorem}\label{thm:B}
Let $Q_B = T + s$ be the def-only pentagon of Example~\textup{\ref{ex:pentagon}},
embedded with vertices $(0,0),(-2,-1),(-3,-2),(-2,-3),(-1,-2)$, and let
\[
\Delta_B=\conv\bigl(Q_B\times\{(1,0)\}\ \cup\ \{e_4,\ (1,1,-1,-1)\bigr\}
\subset \ZZ^4 .
\]
Then:
\begin{enumerate}
\item $\Delta_B$ is reflexive with $7$ vertices and $7$ facets; all $16$ edges
have lattice length one; its $16$ two-dimensional faces are
$F_B = Q_B \times \{(1,0)\}$ and $15$ standard triangles.
\item The generic anticanonical $X_B \subset \PP_{\Delta_B}$ is a Calabi--Yau
threefold whose singular locus consists of exactly $\ell(F_B^\circ)=3$ points,
each with germ $C(Q_B)$.
\item Each singular germ of $X_B$ admits a non-trivial one-parameter
deformation: its reduced miniversal base is a smooth curve.  Nevertheless the
miniversal base has no smoothing component, and $X_B$ admits no smoothing.
\item $(h^{1,1},h^{2,1}) = (9,81)$ for the maximal projective crepant partial
resolution, hence Euler number $-144$.
\end{enumerate}
\end{theorem}

\begin{proof}
(1) The seven facet normals are
\[
\begin{gathered}
(0,0,1,1),\ (0,0,1,-2),\ (-3,6,1,1),\ (6,-3,1,1),\\
(-3,3,-2,1),\ (-1,-1,-4,1),\ (3,-3,-2,1),
\end{gathered}
\]
each at level $1$; the rest is finite verification as in
Theorem~\ref{thm:A}(1), with $F_B^\circ$ again of length $3$.
(2) is Proposition~\ref{prop:planting} plus the face inventory.
(3) is Example~\ref{ex:pentagon} combined with
Theorem~\ref{thm:trichotomy} and Lemma~\ref{lem:localglobal}.
(4) is computed as in Theorem~\ref{thm:A}(4), with
$\ell(\Delta_B)=10$, $\ell(\Delta_B^\circ)=104$.
\end{proof}

\begin{remark}\label{rem:stuck}
By Lemma~\ref{lem:cascade}, the available local deformation at each singular
point of $X_B$ replaces the pentagon cone by a $\frac13(1,1,1)$ point.  Thus,
to the extent that these local deformations are realized by global ones,
$X_B$ moves in a family whose neighbouring members are Calabi--Yau threefolds
with three rigid quotient points --- which then admit no further deformation of
their germs.  In every case the singularities never disappear.  We are not
aware of a previous compact Calabi--Yau example in which all singular points
are locally deformable and yet no smoothing exists; the known non-smoothable
examples (quotient points, and the $\FF_1$-cone of Theorem~\ref{thm:A}) are
locally rigid.
\end{remark}

\subsection{A single non-smoothable point}

The examples above have three singular points because the dual edge of the
planted face has lattice length $3$, and within the two-vertex completion
family this cannot be pushed below $2$: every facet of $\Delta$ containing
the planted face $F$ has normal of the form $u_t = (0,0,1,t)$ with
$t \in \ZZ$ --- these are exactly the integral functionals restricting to $1$
on the affine span of $F$ --- and the two facets through $F$ are the extreme
supporting members of this pencil, so that
$\ell(F^\circ) = t_{\max} - t_{\min}$.  With completing vertices $e_4$ and
$w = (w_1,w_2,w_3,w_4)$ the bound $\ell(F^\circ)\ge 2$ is forced, as
follows.  All vertices of $Q\times\{(1,0)\}$ have $x_3 = 1$ and $x_4 = 0$,
and $e_4$ has $x_3 = 0$, $x_4=1$; for the origin to be interior some
vertex must have negative third coordinate and some vertex negative
fourth coordinate, so
$w_3 \le -1$ and $w_4 \le -1$.  The member $u_t$ supports $\Delta$ if and
only if $\langle u_t, e_4\rangle = t \le 1$ and
$\langle u_t, w\rangle = w_3 + t w_4 \le 1$; since $w_4 < 0$ the second
constraint bounds $t$ from below only, so $t_{\max} = 1$, attained on the
facet through $e_4$.  The facet at the lower extreme must contain a vertex
outside the plane of $F$, necessarily $w$, so
$t_{\min} = (1-w_3)/w_4$, an integer with positive numerator ($\ge 2$) and
negative denominator; hence $t_{\min} \le -1$ and
$\ell(F^\circ) = 1 - t_{\min} \ge 2$.  A third completing vertex in
$\{x_3 = 1\}$ removes the interiority constraint that forced
$w_3 \le -1$ and realizes $t_{\min} = 0$:

\begin{theorem}\label{thm:C}
Let $Q_A$ be embedded as in Theorem~\textup{\ref{thm:A}} and let
\[
\Delta_C=\conv\bigl(Q_A\times\{(1,0)\}\ \cup\ \{e_4,\ (1,1,-1,0),\
(-2,-2,1,-1)\}\bigr)\subset\ZZ^4 .
\]
Then:
\begin{enumerate}
\item $\Delta_C$ is reflexive with $7$ vertices and $10$ facets; all $18$
edges have lattice length one; its $21$ two-dimensional faces are
$F_C = Q_A\times\{(1,0)\}$ and $20$ standard triangles; and the dual edge
$F_C^\circ$, spanned by the negatives of the facet functionals
$(0,0,1,0)$ and $(0,0,1,1)$, has lattice length $1$.
\item The generic anticanonical hypersurface $X_C \subset \PP_{\Delta_C}$ is
a Calabi--Yau threefold whose singular locus is a \emph{single} point, at
which the germ is the anticanonical cone over $\FF_1$.  Every deformation of
$X_C$ over a reduced irreducible base preserves this germ; in particular
$X_C$ admits no smoothing.
\item The maximal projective crepant partial resolution $\widehat X_C$ has
$(h^{1,1},h^{2,1}) = (4,108)$, hence Euler number $-208$.
\end{enumerate}
\end{theorem}

\begin{proof}
The ten facet normals are
\[
\begin{array}{lllll}
(0,0,1,1), & (0,0,1,0), & (-2,0,-3,0), & (-2,0,-3,1), & (-2,4,1,-4),\\[2pt]
(-2,4,1,1), & (0,-2,-3,0), & (0,-2,-3,1), & (4,-2,1,-4), & (4,-2,1,1),
\end{array}
\]
each at level exactly $1$; the face inventory, edge lengths, and the dual
edge $F_C^\circ$ (spanned by the negatives of the first two normals, of
lattice length $1$)
are finite verifications as in Theorem~\ref{thm:A}(1).  Parts (2) and (3)
follow exactly as in Theorem~\ref{thm:A}, with
$\ell(\Delta_C) = 9$, $\ell(\Delta_C^\circ) = 141$ in Batyrev's formulas.
\end{proof}

\begin{remark}
The same pair of completing vertices applied to the $\frac13(1,1,1)$
triangle (embedded with vertices $(0,0),(-2,-1),(-1,-2)$) gives a $6$-vertex
reflexive polytope whose generic anticanonical hypersurface is a Calabi--Yau
threefold with a \emph{single} $\frac13(1,1,1)$ point, with resolution Hodge
numbers $(3,111)$ and $\chi = -216$; the classical
$X_9 \subset \PP(1,1,1,3,3)$ has three such points.
\end{remark}

The rigid examples above leave one configuration outstanding: a single
singular point of the \emph{deformable} type.  Our completion searches
realize $\ell(F^\circ)=1$ for def-only faces only on polytopes carrying
additional singular faces; the example below was instead found by scanning
the Kreuzer--Skarke classification itself ($\Delta_D$ appeared in the scan
of the polytopes with at most $10$ vertices; the scan was subsequently
completed over the whole classification, see \S\ref{sec:questions}).

\begin{theorem}\label{thm:D}
Let
\begin{align*}
\Delta_D=\conv\{&(1,0,0,0),\,(0,1,0,0),\,(0,-1,0,0),\,(0,0,1,0),\,
(-3,1,-1,0),\\
&(0,1,1,0),\,(0,0,0,1),\,(-3,2,1,-1),\,(-2,2,2,-1),\,(1,0,1,1)\}
\subset\ZZ^4 .
\end{align*}
Then:
\begin{enumerate}
\item $\Delta_D$ is reflexive with $10$ vertices and $15$ facets; all $30$
edges have lattice length one; its $35$ two-dimensional faces are $34$
standard triangles and one pentagon $F_D$, which is
$GL_2(\ZZ)$-equivalent to the def-only pentagon $Q_B = T + s$ of
Theorem~\textup{\ref{thm:B}}; and the dual edge $F_D^\circ$, spanned by the
negatives of the facet functionals $(0,1,0,1)$ and $(-1,-1,1,1)$, has
lattice length $1$.
\item The generic anticanonical hypersurface $X_D \subset \PP_{\Delta_D}$ is
a Calabi--Yau threefold whose singular locus is a \emph{single} point, with
germ $C(Q_B)$: the reduced miniversal base of the singularity is a smooth
curve, so the singular point of $X_D$ admits a non-trivial one-parameter
local deformation, and yet $X_D$ admits no smoothing.
\item The maximal projective crepant partial resolution $\widehat X_D$ has
$(h^{1,1},h^{2,1}) = (8,118)$, hence Euler number $-220$.
\end{enumerate}
\end{theorem}

\begin{proof}
All statements in (1) are finite verifications in exact arithmetic as in
Theorem~\ref{thm:A}(1): the fifteen facet normals are
\[
\begin{array}{lllll}
(0,1,0,1), & (-1,-1,1,1), & (1,1,0,0), & (1,1,0,-1), & (1,0,1,-1),\\[2pt]
(1,1,-1,1), & (1,1,-1,-3), & (1,1,-3,1), & (1,1,-3,-5), & (1,-1,1,-1),\\[2pt]
(1,-1,1,-3), & (1,-1,-1,1), & (1,-1,-1,-7), & (1,-1,-5,1), & (1,-1,-5,-11),
\end{array}
\]
each at level exactly $1$; the pentagon face $F_D$ has vertices
$5,7,8,9,10$ in the order listed, its edge multiset in the induced lattice
is $\{(1,0),(3,2),(-1,0),(-3,-1),(0,-1)\}$, and an explicit unimodular map
carries it to $E(Q_B)$; the first two facet normals listed are those of the
two facets through $F_D$, and the dual edge they determine has lattice
length $1$.  Part (2) follows from
Proposition~\ref{prop:planting}, Example~\ref{ex:pentagon},
Theorem~\ref{thm:trichotomy} and Lemma~\ref{lem:localglobal} exactly as in
Theorem~\ref{thm:B}; part (3) is Batyrev's formulas with
$\ell(\Delta_D) = 15$, $\ell(\Delta_D^\circ) = 152$.
\end{proof}

\begin{remark}
Theorems~\ref{thm:C} and~\ref{thm:D} together answer the minimality
question for both non-smoothable types: one rigid point, or one deformable
point, on a compact Calabi--Yau threefold, each realized inside the
Kreuzer--Skarke classification.  The scan that produced $\Delta_D$ found
three $10$-vertex polytopes whose unique \emph{non-smoothable} face is
def-only with dual edge of length one; $\Delta_D$ is the one whose other
faces are all smooth (the other two carry additional ordinary
smoothable singular points).  Over the full classification there are
exactly four polytopes whose generic hypersurface has a \emph{single}
singular point of def-only type --- with $10$, $11$, $11$ and $12$ vertices,
$\Delta_D$ being the minimal one --- and in all four the germ is $C(Q_B)$;
see \S\ref{sec:questions}.
\end{remark}

\section{Further remarks and questions}\label{sec:questions}

\subsection{How common is this?}
The planting search that produced $\Delta_A$ and $\Delta_B$ is the simplest
possible: complete the embedded polygon by two vertices, one of which can be
normalized to $e_4$.  This family realizes $18$ of the $87$ non-smoothable
census classes of Proposition~\ref{prop:census} as two-dimensional faces of
reflexive $4$-polytopes (both examples above use the same completing vertex
$(1,1,-1,-1)$), and the count is stable: enlarging the coordinate box for
the completing vertex from $[-3,3]$ to $[-5,5]$ produces no new classes, and
every class realized this way has at most $6$ interior lattice points.
Completing by \emph{three} vertices instead (as in
Theorem~\ref{thm:C}, with both free vertices ranging over the box
$[-2,2]^4$) raises the count to $32$ of the $87$: all $23$ classes with
$i \le 6$ embed, together with $9$ classes with $7 \le i \le 10$.

A complementary and ultimately decisive approach is to scan the
Kreuzer--Skarke classification itself.  We have carried out the test of
Theorem~\ref{thm:trichotomy} on \emph{all} $473{,}800{,}776$ reflexive
$4$-polytopes \cite{KreuzerSkarke00}: $473{,}800{,}775$ of them through the
per-vertex-count mirror of the database that we used, whose files cover
$5$ to $33$ vertices, and the one remaining polytope separately.
Granting that the mirror's entries are distinct members of the
classification --- its per-file counts and Hodge data were validated
throughout the sweep --- the one polytope it omits is the
classification's unique member with more than $33$ vertices, and we
identify it: it is the $36$-vertex
product of two reflexive hexagons, verified directly (ancillary file
\texttt{missing\_polytope.py}).  That polytope's $48$ two-dimensional faces
are $36$ unit parallelograms and $12$ hexagons of $\mathrm{dP}_6$ type, all
of type \textup{(S)} with dual edge of length one, so it contributes to
none of the counts below.  We find:
\begin{itemize}
\item $8.27\%$ (that is, $39{,}175{,}536$ of them) carry a unit-edge
non-smoothable two-dimensional face, so their generic anticanonical
hypersurfaces are non-smoothable Calabi--Yau threefolds by
Corollary~\ref{cor:global}; the fraction rises from $11.7\%$ at $5$
vertices to a peak of $14.0\%$ at $9$ and then declines steadily with the
vertex count, to below $1\%$ beyond $22$ vertices;
\item def-only faces appear from $7$ vertices on (as they must: a
def-only polygon has at least five edges --- a unit-edge triangle is
indecomposable, and a decomposable unit-edge
quadrilateral splits into two segment pairs and is smoothable --- and a
pentagonal two-face together with the two further vertices needed to span
dimension four already requires seven vertices), $497{,}434$ face
occurrences in all;
\item combining with the synthetic plantings above, exactly $65$ of the
$87$ census classes are realized as two-dimensional faces of reflexive
$4$-polytopes.  The $22$ classes that never occur are all of type
\textup{(D)} and all have $i \ge 11$ interior points, so the occurrence of
a class decays with its size and, beyond a threshold, ceases entirely ---
an obstruction we do not explain;
\item unit-edge non-smoothable faces \emph{outside} the census box also
occur (rigid polygons with an edge coordinate exceeding $2$, with up to
five edges in our samples), so the census stratum is a proper part of the
phenomenon.
\end{itemize}
This answers the first natural quantitative question --- what fraction of
the Batyrev landscape is non-smoothable --- completely, and the second ---
which local models occur --- for the census stratum; the unit-edge
non-smoothable faces outside the census box occur but remain to be
classified.  The single-point examples are similarly constrained: among
all reflexive $4$-polytopes with unit edges whose generic hypersurface
has a \emph{single} singular point, and that point of def-only type, the
point is the pentagon cone $C(Q_B)$ of Theorem~\ref{thm:D} in
\emph{every} case (the four such polytopes on the list all give
$C(Q_B)$); no other def-only germ occurs as the sole singularity of a
Batyrev threefold.  (Def-only germs of other classes do occur as the
unique \emph{non-smoothable} point of a hypersurface that also carries
smoothable singular points; it is only as a lone singularity that
$C(Q_B)$ is forced.)

\begin{question}
The faces with non-unit edges --- whose germs are non-isolated, governed
by Filip's $0$-mutability criterion
\cite{Filip25} rather than by Theorem~\ref{thm:trichotomy} --- deserve their
own count.
\end{question}

\subsection{The smoothable faces and the global defect}
Corollary~\ref{cor:global} uses only the non-smoothable local models.  In the
opposite regime --- all two-dimensional faces of type \textup{(S)} --- local
smoothings exist at every singular point, and the question becomes whether
they can be realized simultaneously by a global deformation; for ordinary
double points this is Friedman's defect condition \cite{Friedman86}, and for
Calabi--Yau threefolds with canonical singularities the local-to-global
obstruction theory of \cite{Gross97,Namikawa02} applies.  A combinatorial
criterion on $\Delta$ for global smoothability of the Batyrev hypersurface,
in the spirit of the Fano-threefold results of \cite{Petracci20,CHP24},
seems within reach and is the natural continuation of this note.

\subsection{Mirror symmetry}
Corti--Filip--Petracci conjecture that the smoothing components of
$\Def C(Q)$ correspond to the $0$-mutable Laurent polynomials with Newton
polygon $Q$ \cite{CFP}, a mirror-theoretic statement.  It would be
interesting to understand what the presence of a type \textup{(R)} or
\textup{(D)} face of $\Delta$ --- and hence the non-smoothability of $X$ ---
means for the mirror family of the Batyrev pair $(\Delta, \Delta^\circ)$.

\begin{thebibliography}{CHP24}

\bibitem[Alt95]{Altmann95}
K.~Altmann,
\emph{Minkowski sums and homogeneous deformations of toric varieties},
T\^ohoku Math.\ J.\ (2) \textbf{47} (1995), no.~2, 151--184.

\bibitem[Alt97]{Altmann97}
K.~Altmann,
\emph{The versal deformation of an isolated toric Gorenstein singularity},
Invent.\ Math.\ \textbf{128} (1997), no.~3, 443--479.

\bibitem[Bat94]{Batyrev94}
V.~V. Batyrev,
\emph{Dual polyhedra and mirror symmetry for Calabi--Yau hypersurfaces in
toric varieties},
J.\ Algebraic Geom.\ \textbf{3} (1994), 493--535.

\bibitem[CFP22]{CFP}
A.~Corti, M.~Filip, and A.~Petracci,
\emph{Mirror symmetry and smoothing Gorenstein toric affine 3-folds},
in \emph{Facets of Algebraic Geometry: A Collection in Honor of William
Fulton's 80th Birthday}, vol.~I, London Math.\ Soc.\ Lecture Note Ser.\
\textbf{472}, Cambridge Univ.\ Press, 2022, pp.~132--163; arXiv:2006.16885.

\bibitem[CHP24]{CHP24}
A.~Corti, P.~Hacking, and A.~Petracci,
\emph{Smoothing Gorenstein toric Fano 3-folds},
arXiv:2412.06500 (2024).

\bibitem[CLS11]{CLS}
D.~A. Cox, J.~B. Little, and H.~K. Schenck,
\emph{Toric Varieties},
Graduate Studies in Mathematics \textbf{124}, Amer.\ Math.\ Soc., 2011.

\bibitem[Fil25]{Filip25}
M.~Filip,
\emph{Laurent polynomials and deformations of non-isolated Gorenstein toric
singularities},
arXiv:2504.04486 (2025).

\bibitem[Fri86]{Friedman86}
R.~Friedman,
\emph{Simultaneous resolution of threefold double points},
Math.\ Ann.\ \textbf{274} (1986), 671--689.

\bibitem[FL25]{FriedmanLaza25}
R.~Friedman and R.~Laza,
\emph{Deformations of Calabi--Yau varieties with isolated log canonical
singularities},
arXiv:2306.03755; Int.\ Math.\ Res.\ Not.\ (2025).

\bibitem[FL24]{FriedmanLazaSigma}
R.~Friedman and R.~Laza,
\emph{Deformations of Calabi--Yau varieties with $k$-liminal singularities},
Forum Math.\ Sigma \textbf{12} (2024), e44.

\bibitem[Gra72]{Grauert72}
H.~Grauert,
\emph{\"Uber die Deformation isolierter Singularit\"aten analytischer Mengen},
Invent.\ Math.\ \textbf{15} (1972), 171--198.

\bibitem[GLS07]{GLS}
G.-M. Greuel, C.~Lossen, and E.~Shustin,
\emph{Introduction to Singularities and Deformations},
Springer Monographs in Mathematics, Springer, 2007.

\bibitem[Gro97]{Gross97}
M.~Gross,
\emph{Deforming Calabi--Yau threefolds},
Math.\ Ann.\ \textbf{308} (1997), 187--220; updated version arXiv:alg-geom/9506022v2 (2025).

\bibitem[Kap09]{Kapustka09}
G.~Kapustka,
\emph{Primitive contractions of Calabi--Yau threefolds II},
J.\ Lond.\ Math.\ Soc.\ (2) \textbf{79} (2009), no.~1, 259--271;
arXiv:0707.2488.

\bibitem[KK09]{KapustkaKapustka09}
G.~Kapustka and M.~Kapustka,
\emph{Primitive contractions of Calabi--Yau threefolds I},
Comm.\ Algebra \textbf{37} (2009), no.~2, 482--502; arXiv:math/0703810.

\bibitem[KS00]{KreuzerSkarke00}
M.~Kreuzer and H.~Skarke,
\emph{Complete classification of reflexive polyhedra in four dimensions},
Adv.\ Theor.\ Math.\ Phys.\ \textbf{4} (2000), 1209--1230.

\bibitem[Nam02]{Namikawa02}
Y.~Namikawa,
\emph{Stratified local moduli of Calabi--Yau threefolds},
Topology \textbf{41} (2002), 1219--1237.

\bibitem[NS95]{NamikawaSteenbrink95}
Y.~Namikawa and J.~H.~M. Steenbrink,
\emph{Global smoothing of Calabi--Yau threefolds},
Invent.\ Math.\ \textbf{122} (1995), 403--419.

\bibitem[Pet20]{Petracci20}
A.~Petracci,
\emph{Some examples of non-smoothable Gorenstein Fano toric threefolds},
Math.\ Z.\ \textbf{295} (2020), 751--760; arXiv:1804.07960.

\bibitem[Sch71]{Schlessinger71}
M.~Schlessinger,
\emph{Rigidity of quotient singularities},
Invent.\ Math.\ \textbf{14} (1971), 17--26.

\end{thebibliography}

\bigskip
\noindent{\sc Bernd J. Wuebben, New York, NY,} \texttt{wuebben@gmail.com}

\end{document}
